\centerline{\bf \Huge Комбинаторный разнобой}

\begin{flushright}
        \scriptsize{\it Меня создало общение с людьми.\\
                        Они создают меня, я создаю их.\\
                        (с) Рей Аянами.}
\end{flushright}

\problem{1}
Эльзяте и Олежке на лбу написали по одной несократимой дроби. Они видят дроби друг друга, но не видят свои дроби. Они должны написать на листе бумаги (каждый на своем листе) конечный набор дробей. Могут ли они действовать так, чтобы хотя бы один из них написал дробь, написанную у него на лбу?

\problem{2}
В классе 33 ученика, а сумма их возрастов составляет 430 лет. Верно ли, что найдутся 20 учащихся, сумма возрастов которых больше 260 лет?

\problem{3}
Несколько прямых, никакие две из которых не параллельны, разрезают плоскость на части. Внутри одной из этих частей отметили точку $A$. Докажите, что точка, лежащая с $A$ по разные стороны от всех данных прямых, существует тогда и только тогда, когда часть, содержащая $A$, неограничена.

\problem{4}
Натуральное число кратно 4. Все его делители выписали в строку в порядке возрастания. Докажите, что в этой строке найдется пара соседей с разностью 2.

\problem{5}
На шахматной доске $n \times n$ расставлены ладьи так, что они бьют все черные клетки. Какое наибольшее возможное количество непобитых белых клеток может быть?

\problem{6.}
Пусть $a_1, a_2, \ldots, a_{99}$ - произвольная перестановка чисел $1,2, \ldots, 99$. Докажите, что среди чисел $\left|a_1-1\right|,\left|a_2-2\right|, \ldots,\left|a_{99}-99\right|$ есть два одинаковых.

\problem{7}
В правильном десятиугольнике проведены все диагонали. Возле каждой вершины и возле каждой точки пересечения диагоналей поставлено число +1 (рассматриваются только сами диагонали, а не их продолжения). Разрешается одновременно изменить все знаки у чисел, стоящих на одной стороне или на одной диагонали. Можно ли с помощью нескольких таких операций изменить все знаки на противоположные?

\problem{8}
Алиса и Боб играют в игру. Игровое поле представляет из себя клетчатую полоску размером $1 \times 2022$. Игроки по очереди выписывают в пустую клетку любую из букв О и Г. Побеждает тот, после чьего хода в трех соседних клетках появятся буквы ОГО. Если все клетки заполнены, а слова ОГО нет, игра заканчивается вничью. Каков будет исход при правильной игре?

\problem{9}
Дана колода карточек. На каждой карточке написано одно из чисел $1,2, \ldots, n$. Известно, что сумма всех чисел на всех карточках равна $k \cdot n!$ для некоторого натурального $k$. Докажите, что можно разложить карточки на $k$ стопок так, чтобы сумма чисел, записанных на карточках в каждой стопке, равнялась бы $n!$.